\documentclass[11pt,letterpaper]{article}

% Use the standard OpenType Latin Modern text fonts with XeLaTeX/Tectonic.
\renewcommand{\encodingdefault}{TU}
\renewcommand{\rmdefault}{lmr}
\renewcommand{\sfdefault}{lmss}
\renewcommand{\ttdefault}{lmtt}
\usepackage[margin=1.05in]{geometry}
\usepackage{amsmath,amssymb,amsthm}
\usepackage[hidelinks]{hyperref}
% Keep large mathematical operator glyphs in the scalable 10-point face.
\DeclareFontShape{OT1}{cmr}{m}{n}{
  <-7>cmr6 <7-8>cmr7 <8-9>cmr8 <9->cmr10
}{}

\hypersetup{
  pdftitle={Fewer than 3.178 to the n stable matchings},
  pdfsubject={An English guide to the complete Lean proof},
  pdfauthor={},
  pdfkeywords={stable matchings, entropy, Lean, formal verification}
}
\newtheorem{theorem}{Theorem}
\newtheorem{lemma}{Lemma}
\newtheorem{proposition}{Proposition}
\theoremstyle{remark}
\newtheorem*{remark}{Remark}
\DeclareMathOperator{\scount}{stableCount}
\newcommand{\E}{\mathbb{E}}
\newcommand{\N}{\mathbb{N}}
\newcommand{\R}{\mathbb{R}}
\newcommand{\cS}{\mathcal{S}}
\newcommand{\code}[1]{\texttt{#1}}
\numberwithin{equation}{section}
\setlength{\emergencystretch}{1em}
\clubpenalty=10000
\widowpenalty=10000
\displaywidowpenalty=10000

\title{Fewer than \(3.178^n\) stable matchings}
\author{\normalsize Proof and formal verification note}
\date{\small 25 September 2026}

\begin{document}
\maketitle

\begin{abstract}
For every positive integer \(n\), a strict, complete stable-marriage
instance with \(n\) men and \(n\) women has fewer than
\((1589/500)^n=3.178^n\) stable matchings. The proof bounds the entropy
of the stable matching family through a random revelation order.
Propagation in both directions of each stable-partner list reduces
the conditional support; a simultaneous-cap remainder supplies an
additional saving. A finite charging argument accounts for exceptional
vertices, and rational logarithm certificates establish the final
constant. This note explains the argument accompanying the complete
Lean proof. The theorem has no certificate hypothesis or restriction
on preference patterns.
\end{abstract}

\section{The theorem}
\label{sec:theorem}

A profile is \emph{strict and complete} if every participant strictly
ranks every participant on the other side. A matching is a bijection
between the two sides. A pair blocks a matching when both participants
strictly prefer each other to their assigned partners; a matching is
\emph{stable} if it has no blocking pair. Let \(\cS(I)\) be the family
of all stable matchings of a profile \(I\), and put
\[
  f(n)=\max_I |\cS(I)|,
\]
where the maximum ranges over strict, complete profiles with \(n\)
participants on each side.

\begin{theorem}\label{thm:main}
For every positive integer \(n\) and every strict, complete profile
\(I\) with \(n\) men and \(n\) women,
\begin{equation}\label{eq:main}
  |\cS(I)| < \left(\frac{1589}{500}\right)^n = 3.178^n.
\end{equation}
In particular, \(f(n)<3.178^n\) for all \(n>0\).
\end{theorem}

The rational base is exact. The result applies to the full stable
matching family, with no bound on \(n\) and no structural hypothesis
on rotations or preferences. The condition \(n>0\) is needed for the
strict inequality: at \(n=0\) there is one empty matching.

The Lean theorem is \code{SmpMax.General.stableCount\_lt\_3178\_pow}.
Its definitions encode preference rows as permutations of ranks, with
lower ranks preferred, and count every stable bijection. Appendix~\ref{app:lean}
gives the exact declaration and verification commands. The complete
machine proof is the accompanying source; this note makes no optimality
or priority claim.

\section{Entropy and a common random gap}
\label{sec:gap}

Assume \(\cS(I)\ne\varnothing\), since otherwise the theorem is
immediate. Choose a reference matching \(\mu\) uniformly from
\(\cS(I)\). For a reveal order \(\pi\), expose the men's partners one
coordinate at a time. Let \(X_m(\mu,\pi)\) be the number of partners
of \(m\) occurring among stable matchings that agree with \(\mu\) at
every already exposed coordinate. The reference itself is compatible,
so \(X_m\ge1\).

\begin{lemma}[Finite revelation inequality]\label{lem:entropy}
For any independently chosen finite random reveal order,
\begin{equation}\label{eq:entropy}
  \log |\cS(I)|
  \le \E_{\mu,\pi}\!\left[\sum_m \log X_m(\mu,\pi)\right].
\end{equation}
\end{lemma}

\begin{proof}
For a fixed order, apply the chain rule for entropy to the uniformly
chosen matching. Each conditional entropy is at most the logarithm of
its conditional support size. Sum over the \(n\) coordinates and then
average over the independently chosen order.
\end{proof}

The formal development establishes the corresponding finite counting
inequality, including the empty-family case. Auxiliary labels introduced
below affect the reveal ranks only; the sum in \eqref{eq:entropy}
always has exactly \(n\) terms.

\subsection{Padded stable-partner lists}

Fix \(\mu\) and \(m\). Order all stable partners of \(m\) from best
to worst. Label a partner \(w\) by her reference owner \(\mu^{-1}(w)\).
The labels are distinct, and the reference partner \(\mu(m)\) has label
\(m\), the pivot. Stable-pair opposition implies that a previously
revealed owner on either side bars compatible partners farther beyond
that owner.

Append 24 auxiliary labels at each end of the list. Choose \(\pi\)
uniformly from the permutations of the \(n\) real labels together with
these 48 auxiliary labels. Every man uses the same auxiliary labels
and the same permutation. Each individual padded list has distinct
labels; independence between men's costs is unnecessary.

Let \(L\) and \(R\) count consecutive later-revealed positions on the
two sides of the pivot, stopping at an earlier label or the padded end.
The common gap has size
\begin{equation}\label{eq:gap}
  K=1+L+R.
\end{equation}
Every compatible stable partner lies in this gap. Finite permutation
counts and a telescoping logarithmic sum give the following uniform
bound, valid independently of the list length.

\begin{lemma}[Finite gap bound]\label{lem:gap}
For a positive cutoff \(N\), define
\begin{equation}\label{eq:BN}
  B_N = \sum_{k=2}^{N} \frac{2}{k+1}\log\frac{k}{k-1}
        +\frac{1}{2N}+\frac{3}{2(N+1)}.
\end{equation}
Then \(\E_\pi[\log K]\le B_N\). At cutoff \(N=512\),
\begin{equation}\label{eq:gap-certificate}
  B_{512}<\log\frac{1665987}{500000}=\log 3.331974.
\end{equation}
\end{lemma}

The tail estimate is finite, and \eqref{eq:gap-certificate} is proved
by an exact rational certificate. The argument needs no identification
of \(B_N\) with an infinite-series limit.

\section{Propagation and directional savings}
\label{sec:directional}

Fix one direction of each stable-partner list. Let \(g(m)\) be the
reference owner of the nearest partner in that direction, or \(m\)
if no such partner exists. Crossing a worse stable partner forces her
reference owner to become worse off, and that change propagates to the
owner's nearest worse-partner owner. The corresponding implication
toward better partners follows from stable-pair opposition in the
original profile.

Thus crossing a partner owned by \(a\) forces \(g(a)\) to change.
If \(g(a)\) has already been revealed, compatibility with \(\mu\)
is contradicted. The crossed partner and every farther partner are
excluded. Both directions are proved for the original preferences.

Let \(l,r\) be the actual left and right compatible-support caps about
the reference partner. They satisfy
\begin{equation}\label{eq:caps}
  0\le l\le L,\qquad 0\le r\le R,\qquad
  1\le X_m\le T:=1+l+r.
\end{equation}
The compatible support need not fill the interval between its extremes;
its cardinality is bounded by the interval length. Measure the two
directional losses against the same gap:
\begin{equation}\label{eq:losses}
  D_L=\log K-\log(1+l+R),\qquad
  D_R=\log K-\log(1+L+r).
\end{equation}
Both losses are nonnegative. Write
\begin{equation}\label{eq:constants}
  a_j=\log\frac{j+2}{j+1},\qquad
  \delta=\frac{467}{20000},\qquad
  \varepsilon=\frac{307}{500000}.
\end{equation}

\subsection{Short sides and noncycles}

If a side has at most one real stable partner, its cap is at most one.
Five padded-gap contributions, telescoped before averaging where they
are nested, yield a directional saving of at least
\begin{equation}\label{eq:alpha}
  \alpha=\frac{a_1}{12}+\frac{a_2}{10}+\frac{a_3}{15}
  \ge 3(\delta+\varepsilon).
\end{equation}
If \(g^2(m)\ne m\), the first partner has a propagation witness outside
the pivot and first neighbor. Three disjoint events, with opposite
gap lengths 0, 1, and 2, give at least
\begin{equation}\label{eq:beta}
  \beta=\frac{a_0}{12}+\frac{a_1}{30}+\frac{a_2}{60}
  \ge 3(\delta+\varepsilon).
\end{equation}
The possible coincidences of that witness with the first three opposite
labels are treated separately; each gives the same certified lower
bound or a larger one.

\subsection{Regular second-step events}

Suppose the chosen side has at least two stable partners and
\(g^2(m)=m\). Let \(b\) own the second partner. In the regular case,
\(u=g(b)\) differs from the pivot, first neighbor, and second neighbor.
If \(u\) is revealed before the pivot, the chosen-side cap is at most one.

For \(j=0,\ldots,23\), require \(u\) and opposite label \(j+1\)
before the pivot, and the first \(j\) opposite labels and first two
chosen-side labels after it. The opposite gap then has length exactly
\(j\), the chosen gap has length at least two, and the loss is at
least \(a_{j+1}\). These events are disjoint for distinct \(j\).

For a witness outside the 24 opposite labels, the probability is
\begin{equation}\label{eq:outside-row}
  p_j=\frac{2}{(j+5)(j+4)(j+3)}.
\end{equation}
If \(u\) is opposite label \(i+1\), where \(0\le i\le23\), it is
\begin{equation}\label{eq:collision-rows}
  p_{ij}=\begin{cases}
    \dfrac{2}{(j+5)(j+4)(j+3)}, & j<i,\\[6pt]
    \dfrac{1}{(j+4)(j+3)}, & j=i,\\[6pt]
    0, & j>i.
  \end{cases}
\end{equation}
Indeed, for \(p\) specified distinct labels before the pivot and \(q\)
specified distinct labels after it, the exact probability is
\begin{equation}\label{eq:order-probability}
  \frac{p!\,q!}{(p+q+1)!}.
\end{equation}
This is a relative-order count, with no independence assumption on
the label events. Designating the outside row by \(i=24\), all 25
rows satisfy the exact certified inequality
\begin{equation}\label{eq:row-bound}
  \sum_{j=0}^{23}p_{ij}a_{j+1}\ge\delta.
\end{equation}

\section{The simultaneous-cap remainder}
\label{sec:joint}

The two directional losses leave a further gain when both caps shrink.
The identity
\begin{equation}\label{eq:product}
  (1+l+R)(1+L+r)-KT=(L-l)(R-r)\ge0
\end{equation}
gives the nonnegative remainder
\begin{align}
  J
  &:=\log\frac{(1+l+R)(1+L+r)}{KT}\notag\\
  &=\log\left(1+\frac{(L-l)(R-r)}{KT}\right)\ge0.
  \label{eq:joint-remainder}
\end{align}
Consequently,
\begin{equation}\label{eq:support-loss}
  \log X_m\le\log K-D_L-D_R-J.
\end{equation}

\begin{lemma}[Simultaneous saving]\label{lem:joint}
If \(L,R\ge2\) and \(l,r\le1\), then \(J\ge\log(16/15)\).
\end{lemma}

\begin{proof}
We have \(T\le3\), and
\(5(L-1)(R-1)\ge1+L+R=K\). Hence
\[
  \frac{(L-l)(R-r)}{KT}
  \ge\frac{(L-1)(R-1)}{3K}\ge\frac1{15}.
\]
Apply the monotonicity of the logarithm in \eqref{eq:joint-remainder}.
The local bound is attained at \(L=R=2\) and \(l=r=1\).
\end{proof}

For regular second-step witnesses in both directions, require the
witnesses before the pivot and the four nearest side labels after it.
Distinct witnesses avoiding these four labels give probability
\(2!4!/7!=1/105\); coincident witnesses give \(1!4!/6!=1/30\).
In either case the expected joint saving is at least \(\varepsilon\),
because
\begin{equation}\label{eq:epsilon}
  \frac{1}{105}\log\frac{16}{15}\ge\varepsilon.
\end{equation}

A witness may instead coincide with one of the first two opposite
labels. This prevents the simultaneous event but increases its own
directional saving. The two relevant lower bounds are
\begin{equation}\label{eq:cross-collision}
  \frac{a_1}{12}\ge\delta+\varepsilon,
  \qquad
  \frac{a_1}{30}+\frac{a_2}{20}\ge\delta+\varepsilon.
\end{equation}
The other direction retains its \(\delta\) saving. Thus every regular
cross-direction witness configuration gives combined saving at least
\(2\delta+\varepsilon\), including these collisions.

\section{Exceptional vertices and charging}
\label{sec:charging}

The joint bonus is an aggregate bound, since it cannot be assigned to
each man separately. For each direction, define the noncycle set
\[
  B=\{m:g^2(m)\ne m\}.
\]
Let \(D\) contain the exceptional men whose chosen side has at least
two partners, with \(g^2(m)=m\) and \(g(b)\in\{m,g(m)\}\), where
\(b\) owns the second partner.

\begin{lemma}[Two-to-one charging]\label{lem:charging}
For each direction, \(|D|\le2|B|\).
\end{lemma}

\begin{proof}
Charge an exceptional man \(m\) to its second-partner owner \(b\).
The neighbor-owner activity lemma gives \(g(b)\ne b\), and the
exceptional relations imply \(g^2(b)\ne b\), so \(b\in B\).
Conversely, a man charging a fixed \(b\) must be either \(g(b)\) or
\(g^2(b)\). Each fiber therefore has at most two elements.
\end{proof}

Each directional mean loss is nonnegative, at least \(\delta\) outside
\(D\), and at least \(3(\delta+\varepsilon)\) on \(B\). Short sides
use \eqref{eq:alpha}; regular sides use the event rows of
Section~\ref{sec:directional}. A man regular on both sides has the
combined gain from Section~\ref{sec:joint}. If either side is short,
its stronger saving already covers that combined target.

Put \(S_m=\E_\pi[D_L+D_R+J]\). Write \(b_L,b_R\) and \(d_L,d_R\)
for the indicators, at \(m\), of the two noncycle and exceptional
sets. The resulting pointwise accounting inequality is
\begin{equation}\label{eq:accounting}
  S_m\ge2\delta+\varepsilon
       +2(\delta+\varepsilon)(b_L+b_R)
       -(\delta+\varepsilon)(d_L+d_R).
\end{equation}
Summing over men and applying Lemma~\ref{lem:charging} separately
in each direction cancels the possible deficits:
\begin{equation}\label{eq:total-saving}
  \sum_m S_m\ge n(2\delta+\varepsilon),
  \qquad
  2\delta+\varepsilon=\frac{23657}{500000}=0.047314.
\end{equation}
This estimate includes the cost of every exceptional vertex.

\section{Exact constants and the final bound}
\label{sec:assembly}

Combining the gap estimate, \eqref{eq:support-loss}, and
\eqref{eq:total-saving} gives, for every stable reference \(\mu\),
\begin{equation}\label{eq:cost-bound}
  \E_\pi\!\left[\sum_m\log X_m(\mu,\pi)\right]
  \le n(B_{512}-2\delta-\varepsilon).
\end{equation}
The remaining numerical comparison is exact:
\begin{equation}\label{eq:final-certificate}
  B_{512}-2\delta-\varepsilon<\log\frac{1589}{500}.
\end{equation}

To certify the event bounds and this comparison, the development uses
finite positive partial sums for the logarithm. For \(0\le x<1\)
and an integer \(t\ge1\),
\begin{equation}\label{eq:log-series}
  \log\frac{1+x}{1-x}
  =2\sum_{k=0}^{t-1}\frac{x^{2k+1}}{2k+1}+R_t(x),
  \qquad
  0\le R_t(x)\le
  \frac{2x^{2t+1}}{(2t+1)(1-x^2)}.
\end{equation}
For \(a_j\), take \(x=1/(2j+3)\); six positive terms give the
rational floors used in the 25 event rows. For the final comparison,
take
\begin{equation}\label{eq:ratio}
  x=\frac{76987}{3254987},
  \qquad
  \frac{1+x}{1-x}
  =\frac{1665987/500000}{1589/500}.
\end{equation}
Four terms and the rational remainder bound certify that the logarithm
of this ratio is less than \(2\delta+\varepsilon\). Together with
\eqref{eq:gap-certificate}, this proves \eqref{eq:final-certificate}.
All comparisons are proved in Lean by exact arithmetic.

\begin{proof}[Proof of Theorem~\ref{thm:main}]
The empty stable family is immediate. Otherwise, average
\eqref{eq:cost-bound} over \(\mu\) and use
Lemma~\ref{lem:entropy}. Since \(n>0\), the strict inequality
\eqref{eq:final-certificate} and monotonicity of the exponential give
\begin{align*}
  |\cS(I)|
  &\le\exp\!\bigl(n(B_{512}-2\delta-\varepsilon)\bigr)\\
  &<\exp\!\left(n\log\frac{1589}{500}\right)\\
  &=\left(\frac{1589}{500}\right)^n.
\end{align*}
Taking the maximum over profiles completes the proof.
\end{proof}

\appendix
\section{The Lean statement and verification}
\label{app:lean}

The exact public declaration in
\code{JointEntropyUpperBound.lean} is:
\begin{quote}\small\ttfamily
theorem stableCount\_lt\_3178\_pow \{n : \(\N\)\}\\
\hspace*{1em}(hn : 0 < n) (I : Profile n) :\\
\hspace*{1em}(stableCount I : \(\R\)) < (1589 / 500 : \(\R\)) \textasciicircum\ n
\end{quote}
It is in namespace \code{SmpMax.General}. The count is coerced to
\(\R\); \(1589/500\) is a rational expression in \(\R\), not a
floating-point approximation. The exact statement check and definition
printout are in \code{lean/Check3178.lean}.

The final theorem's axiom dependencies are exactly
\code{propext}, \code{Classical.choice}, and \code{Quot.sound}.
There is no \code{sorry}, \code{admit}, new axiom, or
\code{native\_decide} in the packaged general proof sources. Numerical
validation scripts and external SAT certificates are not assumptions
of the theorem.

The bundle contains 56 unchanged general proof modules. The target's
transitive local dependency closure comprises 45 modules; the remaining
modules retain earlier completed bounds. Lean is pinned to version
4.33.1, and the exact Mathlib and transitive dependency revisions are
recorded in \code{lean/lake-manifest.json}.

\begin{samepage}
From the bundle root, with Python 3.9 or newer and the Lean toolchain
manager \code{elan} installed, run:
\begin{quote}\small
\begin{verbatim}
python3 tools/verify.py --fetch
python3 tools/verify.py --kernel all
\end{verbatim}
\end{quote}
\end{samepage}
The first command obtains the pinned dependencies, builds all 56 source
modules, checks the exact declaration, audits 112 general statements,
and independently replays the final module. The second also replays
all 45 local target-dependency modules. Omit \code{--fetch} when
dependencies are already present. For a file-identity check without Lean:
\begin{quote}\small
\begin{verbatim}
python3 tools/verify.py --sources-only
\end{verbatim}
\end{quote}

The retained verification was performed on 24 September 2026: all 56
modules rebuilt, all 112 audits passed, and all 45 modules replayed
successfully. It reused the pinned dependency cache; the download step
was not repeated. This 25 September edition changes the presentation
only. The Lean sources, project configuration, verifier, and formal
verification logs are byte-identical to that checked bundle. The logs
and command results are in \code{evidence/verification.json} and its
adjacent files.

\section{Guide to the formal development}
\label{app:map}

All module names below are relative to \code{lean/SmpMax/General/}
and have the suffix \code{.lean}. The complete source list and dependency
order are in \code{SOURCE\_MAP.md} and \code{source-index.json}.

\begin{description}
  \raggedright
  \item[Definitions and revelation.]
  \code{Definitions}, \code{MatchingRevelation}, and
  \code{ExtendedRevelation} define the full stable family and establish
  the finite revelation inequality in Section~\ref{sec:gap}.

  \item[Common-gap bound.]
  \code{PartnerIntervals}, \code{PermutationCounting},
  \code{LogTail}, and \code{RefinedLogBound} prove the finite gap
  estimate. \code{RefinedLogCertificate} checks its exact constant;
  \code{LongPaddedOwners} and \code{LongPaddedSupport} supply the
  longer padded-list interfaces.

  \item[Propagation and event rows.]
  \code{SuccessorPropagation}, \code{BidirectionalPropagation},
  \code{SupportCaps}, and \code{PartnerFrame} connect exclusions to
  actual compatible supports. \code{LongDirectionalSaving},
  \code{LongEventRows}, and \code{LongerLogCertificate} establish
  the directional events and all 25 certified rows.

  \item[Joint savings and charging.]
  \code{JointGapSaving}, \code{JointGapGeometry}, and
  \code{JointDirectionalSaving} cover the simultaneous remainder
  and witness collisions. \code{SuccessorCharging},
  \code{JointCharging}, and \code{JointMatchingSaving} establish
  the global saving, including exceptional vertices.

  \item[Exact final comparison.]
  \code{JointLogCertificate} proves
  \eqref{eq:final-certificate}; \code{JointEntropyUpperBound}
  assembles the support estimates and entropy inequality into
  Theorem~\ref{thm:main}.
\end{description}

\end{document}
